\documentclass[11pt,a4paper]{article}
\usepackage[UTF8]{ctex}
\usepackage{amsmath, amssymb, amsthm}
\usepackage{geometry}
\geometry{
    twoside,      % 开启双面排版模式，使奇偶页边距镜像对称
    inner=1.8cm,  % 控制距物理纸张外边缘的距离（调大，此数值对应左页的左侧、右页的右侧）
    outer=0.5cm,  % 控制物理纸张两页中间缝隙的一半（调小，中间总缝隙 = 0.5 + 0.5 = 1cm）
    top=1.4cm,    
    bottom=1.5cm
}


% --- 调整列表环境间距 ---
\usepackage{enumitem}
\setlist[enumerate]{itemsep=0pt, parsep=0pt, topsep=2pt, partopsep=0pt}
\setlist[itemize]{itemsep=0pt, parsep=0pt, topsep=2pt, partopsep=0pt}
% --- 定理类环境（共享计数器，按 4.x 格式编号）---
\makeatletter
\def\thm@space@setup{%
  \thm@preskip=4pt plus 2pt minus 2pt
  \thm@postskip=\thm@preskip
}
\makeatother
\theoremstyle{plain}
\newtheorem{thm}{定理}
\newtheorem*{thm*}{定理}
\theoremstyle{definition}
\newtheorem{defn}[thm]{定义}
\newtheorem{lem}[thm]{引理}
\newtheorem{cor}[thm]{推论}
\newtheorem{prop}[thm]{命题}
\renewcommand{\thethm}{4.\arabic{thm}}
\renewcommand{\proofname}{证明}

%================ 缩小公式上下及多行公式间距 ================
\AtBeginDocument{
  \abovedisplayskip=4pt plus 1pt minus 1.5pt
  \belowdisplayskip=4pt plus 1pt minus 1.5pt
  \abovedisplayshortskip=0pt plus 3pt
  \belowdisplayshortskip=3pt plus 1pt minus 1pt 
}
\setlength{\jot}{1pt} % 缩小 align 等环境中的行间距
\pagestyle{empty} % 去除页眉页脚（包括页码）

\begin{document}

% ========================================== %
% 一、引入 Artin + 上节衔接
% ========================================== %

\noindent \textbf{Emil Artin} 主要贡献包括类域论的公理化、Artin L-函数、Artin 环与辫群.Wedderburn证明了有限除环都是域.关于本节的定理，Wedderburn 于 1908 年证明了有限维代数的结构定理，Artin 于 1927 年将其推广到满足降链条件的环——这就是今天要讲的 Wedderburn-Artin 定理的历史脉络.课本所讲是有限维版本，而定理的最终形态属于 Artin 环理论.\par



\noindent \textbf{过渡：} 既然每个有限维代数的左正则模都有合成列，一个自然的问题是：什么情况下这些合成因子都是单模，且左正则模本身就是单子模的直和？这就是今天要研究的半单代数.

% ========================================== %
% 二、引子——群表示的两个问题
% ========================================== %


\noindent 让我们先回顾群表示论中的两个基本事实.设 $G$ 为有限群，$F$ 为特征不整除 $|G|$ 的分裂域，那么\par
\noindent \textbf{1：} $G$ 的不可约表示个数恰好等于 $G$ 的共轭类个数.\par
\noindent \textbf{2：} 各不可约表示的维数平方和等于群的阶，即 $\sum n_i^2 = |G|$.\par
%现在我们想要将上述结论推广到更一般的代数框架中，也就是我们要研究一个代数有多少个单模，以及这些单模之间的维数关系.\par

% ========================================== %
% 三、半单代数定义
% ========================================== %

\par
由于上述讨论都基于完全可与表示，因此类似地，我们需要先定义半单代数的概念.

\setcounter{thm}{0}

\begin{defn}
设 $A$ 是有限维 $F$-代数.如果左正则 $A$-模 ${}_A A$ 是半单模，即 ${}_A A$ 是其单子模的直和，则称 $A$ 是半单代数.
\end{defn}

如果 $A$ 的非零左理想$I$ 不再真包含 $A$ 的非零左理想，则称其称为极小左理想.介绍极小左理想的原因是我们需要下列引理.

\begin{lem}
设 $A$ 是 $F$-代数.则左正则模 ${}_A A$ 的子模是$A$ 的左理想.${}_A A$ 的单子模是$A$ 的极小左理想. \par
特别地，$A$ 是半单代数当且仅当 $A$ 是其极小左理想的和.
\end{lem}
\begin{proof}
${}_A A$ 作为 $F$-向量空间就是 $A$ 自身.其 $A$-子模 $M \subseteq {}_A A$ 需满足：
\begin{enumerate}
    \item $M$ 是加法子群，
    \item 对任意 $a \in A$, $m \in M$，有 $a m \in M$.
\end{enumerate}
条件 (i) 和 (ii) 恰好是一个左理想的定义.\par
进一步，若$M$是${}_A A$的单子模，则$M$没有非零真子模.对应地，作为左理想，$M$不真包含任何非零左理想，故为极小左理想.\par
由于$M$是半单模当且仅当它是若干单子模的和，因此$A$是半单代数当且仅当$A$是其极小左理想的和.(书上引理2.5第一条).
\end{proof}

\noindent\textbf{例：}可除代数 $A$ 是半单代数.

可除代数只有两个左理想：$\{0\}$ 和 $A$ 自身.由上述引理即可得$A$ 是半单代数.\par
这同时说明可除代数不仅是半单的，更是单代数—.这就是 Wedderburn-Artin 定理中将可除代数作为基本积木块的原因.


\noindent\textbf{例：}若 $G$ 是有限群且 $\mathrm{char}F \nmid |G|$，由 Maschke 定理知群代数 $FG$ 是半单代数.

% ========================================== %
% 四、Wedderburn-Artin 定理
% ========================================== %

\bigskip
在进入主定理之前,需要两个预备引理.第一个引理的想法是，不可约表示是一个完全可约表示，单模是半单模，那么单代数是否确实是是一个半单代数.

\begin{lem}
单代数是半单代数.
\end{lem}

\begin{proof}
设 $A$ 是单代数.由引理~4.2，只需证 $A$ 是其极小左理想的和.令 $M$ 为 $A$ 的所有极小左理想之和.因 $A$ 是有限维代数，极小左理想存在，故 $M \neq 0$.

任取极小左理想 $S$ 和 $a \in A$.定义 $\varphi_a: S \to A, s \mapsto sa$.$\varphi_a$ 是左 $A$-模同态，其像 $Sa$ 是 $A$ 的子模.由同态基本定理，$Sa \cong S / \ker\varphi_a$.因 $S$ 为单模，$\ker\varphi_a = 0$ 或 $\ker\varphi_a = S$.前者给出 $Sa \cong S$，故 $Sa$ 也是极小左理想，从而 $Sa \subseteq M$；后者给出 $Sa = 0 \subseteq M$.

这对每个极小左理想 $S$ 和每个 $a \in A$ 成立，故 $M \cdot A \subseteq M$，即 $M$ 是 $A$ 的右理想.$M$ 作为左理想之和自然是左理想.因此 $M$ 是 $A$ 的非零双边理想.由 $A$ 的单性，$M = A$，即 $A$ 是其极小左理想的和，故为半单代数.
\end{proof}


Wedderburn-Artin 证明中还会出现 $(\mathrm{End}_A(A))^{\mathrm{op}}$ 这一项，因此还需介绍反代数的概念.


\noindent \textbf{反代数的定义与例子.} 设 $A$ 是 $F$-代数，其乘法记为 $a \cdot b$.定义 $A$ 的\textbf{反代数} $A^{\mathrm{op}}$ 如下：
\begin{itemize}
    \item 作为 $F$-向量空间，$A^{\mathrm{op}} = A$（集合完全相同）；
    \item 乘法重新定义为 $a \circ b := b \cdot a$（把原乘法的顺序反过来）.
\end{itemize}
容易验证 $A^{\mathrm{op}}$ 仍为 $F$-代数，且 $(A^{\mathrm{op}})^{\mathrm{op}} = A$.

下面给出关键的结论：$\mathrm{End}_A({}_A A) \cong A^{\mathrm{op}},\mathrm{End}_A({}_A A)^{\mathrm{op}}\cong A$.

\noindent \emph{证明：} 定义映射
\begin{align*}
\Phi: A^{\mathrm{op}} &\to \mathrm{End}_A({}_A A) \\
a &\mapsto \Phi(a):\begin{aligned}[t]
{}_A A &\to {}_A A \\
x &\mapsto x a
\end{aligned}
\end{align*}
\begin{itemize}
    \item \textbf{良定性：} 对任意 $b \in A$，$\Phi(a)(b x) = (b x)a = b(x a) = b \cdot \Phi(a)(x)$，故 $\Phi(a)$ 是左 $A$-模同态.
    \item \textbf{保持乘法：} 在 $A^{\mathrm{op}}$ 中乘法为 $a \circ b = ba$，则
    \[
    \Phi(a \circ b)(x) = \Phi(ba)(x) = x(ba) = (x b)a = \Phi(a)(x b) = \Phi(a)(\Phi(b)(x)) = (\Phi(a) \circ \Phi(b))(x),
    \]
    即 $\Phi(a \circ b) = \Phi(a) \circ \Phi(b)$.
    \item \textbf{单射：} 若 $\Phi(a) = 0$，取$x=1$，则 $1 \cdot a = a = 0$.
    \item \textbf{满射：} 任取 $f \in \mathrm{End}_A({}_A A)$，令 $a = f(1)$.则对任意 $x \in A$，
    \[
    f(x) = f(x \cdot 1) = x \cdot f(1) = x a = \Phi(a)(x),
    \]
    故 $f = \Phi(a)$.
\end{itemize}

因此 $\Phi$ 是 $F$-代数同构，$A^{\mathrm{op}} \cong \mathrm{End}_A(A)$.\par
再取反即得 $A \cong (\mathrm{End}_A(A))^{\mathrm{op}}$


\medskip
然后可以介绍下面有关矩阵代数的引理.

\begin{lem}
设 $D$ 是可除 $F$-代数，$(M_n(D))^{\mathrm{op}}$ 是 $D$ 上 $n$ 阶全矩阵代数 $M_n(D)$ 的反代数.则有 $F$-代数同构 $(M_n(D))^{\mathrm{op}} \cong M_n(D^{\mathrm{op}})$，其中 $D^{\mathrm{op}}$ 是 $D$ 的反代数.
\end{lem}



\begin{proof}
令 $\pi : (M_n(D))^{\mathrm{op}} \to M_n(D^{\mathrm{op}})$ 将 $D$ 上 $n$ 阶矩阵 $X$ 送到 $X$ 的转置 $X^{\mathrm{T}}$.则 $\pi$ 是 $F$-线性映射.设 $X=(x_{ij})$ 和 $Y=(y_{ij})$ 均是 $n$ 阶矩阵.分别用 $X \circ Y$ 和 $X * Y$ 表示 $(M_n(D))^{\mathrm{op}}$ 和 $M_n(D^{\mathrm{op}})$ 中的乘法.设 $x, y \in D$.用 $x \circ y$ 表示 $D^{\mathrm{op}}$ 中的乘法.则（注意：若 $D$ 不可换，则在 $M_n(D)$ 中 $(YX)^{\mathrm{T}}$ 未必等于 $X^{\mathrm{T}} \cdot Y^{\mathrm{T}}$）

$$
\begin{aligned}
(\pi(X \circ Y))_{ij} &= (\pi(YX))_{ij} = ((YX)^{\mathrm{T}})_{ij} \\
&= (YX)_{ji} = \sum_{k=1}^n y_{jk} x_{ki} \\
&= \sum_{k=1}^n x_{ki} \circ y_{jk} = (X^{\mathrm{T}} * Y^{\mathrm{T}})_{ij} \\
&= (\pi(X) * \pi(Y))_{ij},
\end{aligned}
$$

即 $\pi(X \circ Y) = \pi(X) * \pi(Y)$.从而 $\pi$ 是 $F$-代数同构.
\end{proof}

\medskip
有了上述技术准备，现在转向核心定理.artin-wedderburn定理从四个角度刻画半单代数：
(i) 是内部的（左正则模半单）；（ii）是外部的（所有模半单）；
(iii) 是结构的（矩阵代数直积）；（iv）是分类的（单代数直积）.
其中 (ii) $\Rightarrow$ (iii) 是证明的核心.

下述定理是 Wedderburn-Artin 关于半单代数结构的主要结果.

\begin{thm}[Wedderburn-Artin]
下述命题等价：
\begin{itemize}
    \item[(i)] $A$ 是半单 $F$-代数.
    \item[(ii)] 任一左 $A$-模均是半单模.
    \item[(iii)] $A$ 同构于可除 $F$-代数上全矩阵代数的直积.特别地，若 $F$ 是代数闭域，$A$ 同构于 $F$ 上全矩阵代数的直积.
    \item[(iv)] $A$ 同构于单代数的直积.
\end{itemize}
\end{thm}

\begin{proof}
\noindent \textbf{(i) $\Rightarrow$ (ii)：} 设 $A$ 是半单代数，$M$ 左 $A$-模.取 $M$ 的一组 $F$-基 $m_1, \dots, m_n$.定义满 $A$-模同态
\[
\varphi: A^n \longrightarrow M, \quad (a_1,\dots,a_n) \mapsto \sum_{i=1}^n a_i m_i.
\]
则 $M \cong A^n / \ker\varphi$.由于 $A$ 半单，左正则模 ${}_A A$ 是半单模，故 $A^n$ 作为半单模的直和仍半单.$M$ 是半单模 $A^n$ 的商模，由 2.5 引理 2 知商模保持半单性，故 $M$ 半单.

\medskip
\noindent \textbf{(ii) $\Rightarrow$ (iii)：} 设任一左 $A$-模均是半单模.则左正则模 ${}_A A$ 是半单模.故 $A = n_1S_1 \oplus \dots \oplus n_tS_t$，其中 $S_1, \dots, S_t$ 是 ${}_A A$ 的互不同构的单子模.令 $n_iS_i = L_i$.由于$n_i$都有限，因此可以将其提出，再由Schur 引理知
$$
\mathrm{Hom}_A(L_i, L_j) \cong n_i n_j \mathrm{Hom}_A(S_i, S_j) = 0, \ i \neq j.
$$

再由推论 2.12 知
$$
\mathrm{End}_A(L_i) = \mathrm{End}_A(n_iS_i) \cong M_{n_i}(D_i),
$$
其中 $D_i = \mathrm{End}_A(S_i)$ 是 $F$-上可除代数.

由引理 2.11 和上述引理（反代数同构），有
$$
\begin{aligned}
A &\cong (\mathrm{End}_A(A))^{\mathrm{op}} = (\mathrm{End}_A(L_1 \oplus \dots \oplus L_t))^{\mathrm{op}} \\
&\cong (\mathrm{End}_A(L_1) \times \dots \times \mathrm{End}_A(L_t))^{\mathrm{op}} \\
&\cong (M_{n_1}(D_1) \times \dots \times M_{n_t}(D_t))^{\mathrm{op}} \\
&\cong (M_{n_1}(D_1))^{\mathrm{op}} \times \dots \times (M_{n_t}(D_t))^{\mathrm{op}} \\
&\cong M_{n_1}(D_1^{\mathrm{op}}) \times \dots \times M_{n_t}(D_t^{\mathrm{op}}).
\end{aligned}
$$

因为 $D_i$ 是可除代数，故 $D_i^{\mathrm{op}}$ 也是可除代数，从而 $A$ 同构于可除代数上全矩阵代数的直积.

\medskip
\noindent \textbf{(iii) $\Rightarrow$ (iv)：}只需证明每个直积因子 $M_{n}(D)$（$D$ 可除）是单代数，则 $A$ 自然是单代数的直积.这个结论实际上在3.1的作业第2题已经证明过了.

设 $I \neq 0$ 是 $M_n(D)$ 的双边理想.取 $0 \neq X = (x_{ij}) \in I$，设 $x_{pq} \neq 0$.令 $e_{ij}$ 表示 $(i,j)$ 处为 $1$、其余为 $0$ 的矩阵基.对任意 $i,j$，
\[
e_{ip} X e_{qj} = e_{ip}\Bigl(\sum_{k,\ell} x_{k\ell} e_{k\ell}\Bigr) e_{qj}
= e_{ip} \bigl(x_{pq} e_{pq}\bigr) e_{qj}
= x_{pq} e_{ij} \in I.
\]
由于 $D$ 可除，$x_{pq}$ 可逆，故 $e_{ij} = x_{pq}^{-1} (x_{pq} e_{ij}) \in I$.所以 $I$ 包含所有矩阵基，从而 $I = M_n(D)$.于是 $M_n(D)$ 无非平凡双边理想，即为单代数.

\medskip
\noindent \textbf{(iv) $\Rightarrow$ (i)：} 设 $A = A_1 \times \dots \times A_n$，$A_i$ 均为单代数.由引理~4.3，每个 $A_i$ 是半单代数，即 $A_i$ 是其极小左理想（单 $A_i$-模）的和.

将每个 $A_i$ 与子代数 $\{0\}\times\cdots\times A_i\times\cdots\times\{0\} \subseteq A$ 等同，则 $A_i$ 是 $A$ 的理想，且 $A = A_1 \oplus \dots \oplus A_n$.任一左 $A_i$-模 $M$ 可自然地视为左 $A$-模：对 $a = (a_1,\dots,a_n) \in A$ 和 $m \in M$，定义
\[
a \cdot m := a_i m.
\]
在此作用下，$M$ 的 $A$-子模与 $A_i$-子模一致（因为对 $j \neq i$，$a_j$ 分量不参与作用）.从而单 $A_i$-模提升后仍为单 $A$-模.故每个 $A_i$ 是单 $A$-模的和.

而左正则 $A$-模 ${}_A A = A_1 \oplus \dots \oplus A_n$是$A_i$ 的直和，因此 ${}_A A$ 是单 $A$-模的和，即 $A$ 是半单代数.
\end{proof}

% ========================================== %
% 五、推论展开
% ========================================== %

\noindent Wedderburn-Artin 定理告诉我们半单代数就是可除代数上矩阵代数的直积.现在从这个核心定理出发，逐步展开一系列推论：第一，直积因子——单代数——如何刻画；第二，每个直积因子上有哪些单模；第三，它们一共贡献了多少个互不同构的单模.

\begin{cor}
$A$ 是单代数当且仅当 $A$ 同构于可除代数上的某一全矩阵代数.

特别地，若 $F$ 是代数闭域，则 $A$ 是单代数当且仅当 $A$ 同构于 $F$ 上的某一全矩阵代数.
\end{cor}

\begin{proof}
设 $A$ 是单代数.由引理~4.3 知 $A$ 是半单代数.从而由 Wedderburn-Artin 定理知 $A \cong M_{n_1}(D_1) \times \dots \times M_{n_t}(D_t)$，其中 $D_i$ 均为可除代数.因为每个直积因子 $M_{n_i}(D_i)$ 均为 $A$ 的理想，由 $A$ 的单性即知 $t=1$.
\end{proof}


\noindent 设 $A$ 是半单代数.根据 Wedderburn-Artin 定理，$A$ 就是一些 $M_{n_i}(D_i)$ 的直积.要弄清 $A$ 上所有表示，关键在于两点：单模归属于哪个直积因子，以及每个 $M_n(D)$ 上到底有哪些单模.下述定理完整回答了这两个问题.

\begin{thm}
\begin{itemize}
    \item[(i)] 设 $A = A_1 \times \dots \times A_n$ 是半单代数，其中每个 $A_i$ 均为单代数，$S$ 是任一单 $A$-模.则存在唯一的 $A_i$ 使得 $S$ 是单 $A_i$-模且 $A_j S = 0, \forall j \neq i$.称 $A_i$ 为 $S$ 相应的单因子.从而，半单代数上单模是其某一个单因子上的单模.

    \item[(ii)] 设 $D$ 是有限维可除 $F$-代数.则在同构意义下 $A = M_n(D)$ 有唯一的单模 $V$；且正则 $A$-模同构于 $nV$，$\dim_F V = n[D:F]$；$\mathrm{End}_A(V) \cong D^{\mathrm{op}}$.

    \item[(iii)] 设 $D$ 是有限维可除 $F$-代数，$(V, \rho)$ 是 $A = M_n(D)$ 的不可约表示，其中 $\rho : A \to \mathrm{End}_F(V)$ 是相应的代数同态.则 $\rho(A) \cong M_n(D)$.若 $\mathrm{Hom}_A(V, V) \cong F$，则 $\rho(A) = M_n(F)$.此时 $\dim_F V = n$；并且存在 $V$ 的一组 $F$-基 $B$ 使得 $\rho(a)$ 在 $B$ 下的矩阵就是 $a, \forall a \in A$.
\end{itemize}
\end{thm}

\begin{proof}
\noindent \textbf{(i)} 设 $e_i$ 是 $A_i$ 的单位元，$i = 1, \dots, n$.由于 $A = A_1 \times \dots \times A_n$ 是代数直积，$e_i$ 是中心正交幂等元：$e_i e_j = 0\;(i \neq j)$，$e_i^2 = e_i$，且 $1 = e_1 + \dots + e_n$.\par
\noindent\fbox{$e_j S$ 是 $A$-子模} 对任意 $a \in A$ 和 $s \in S$，因 $e_j$ 在 $A$ 的中心里，有
\[
a \cdot (e_j s) = (a e_j) \cdot s = (e_j a) \cdot s = e_j \cdot (a s) \in e_j S.
\]
故 $e_j S$ 对 $A$ 的作用封闭，是 $S$ 的 $A$-子模.\par
\noindent \fbox{恰有一个 $e_i S = S$，其余为 $0$} $S$ 是单模，故每个子模 $e_j S$ 只能是 $0$ 或 $S$.又
\[
S = 1 \cdot S = (e_1 + \dots + e_n) \cdot S = e_1 S + \dots + e_n S,
\]
而 $S \neq 0$，故至少存在一个 $i$ 使 $e_i S \neq 0$，进而 $e_i S = S$.

若另有 $j \neq i$ 也使 $e_j S = S$，则
\[
e_j S = e_j(e_i S) = (e_j e_i) S = 0 \cdot S = 0,
\]
与 $e_j S = S \neq 0$ 矛盾.因此 $e_j S = 0$ 对所有 $j \neq i$ 成立.\par
\noindent\fbox{$S$ 作为 $A_i$-模是单模} 由上，
\[
A_i S = (A e_i) S = A(e_i S) = A S = S, \qquad
A_j S = (A e_j) S = A(e_j S) = A \cdot 0 = 0 \;(j \neq i).
\]
现设 $N \subseteq S$ 是任一 $A_i$-子模.则对任意 $a = (a_1,\dots,a_n) \in A$，
\[
a N = (a_1 + \dots + a_n) N = a_i N \subseteq N \quad (\text{因 } a_j N \subseteq A_j S = 0 \text{ 对 } j \neq i).
\]
故 $N$ 也是 $A$-子模.由 $S$ 的单性，$N = 0$ 或 $N = S$.因此作为 $A_i$-模，$S$ 没有非平凡子模，即为单模.\par
\noindent \textbf{(ii)} 定义$V$是$D$上的$n$为列向量空间
$$
V := \left\{ \begin{pmatrix} x_1 \\ \vdots \\ x_n \end{pmatrix} \ \middle| \ x_1, \dots, x_n \in D \right\}.
$$
则 $V$ 自然地作成左 $M_n(D)$-模.\par
\noindent\fbox{$V$是单的} 设 $U$ 是 $V$ 的非零子模.任取非零向量$a = \begin{pmatrix} a_1 & \cdots & a_n \end{pmatrix}^T \in U$，其中某一 $a_j \neq 0$，总有
$$
\begin{pmatrix}
0 & \dots & 0 & x_1 a_j^{-1} & 0 & \dots & 0 \\
\vdots & & \vdots & \vdots & \vdots & & \vdots \\
0 & \dots & 0 & x_n a_j^{-1} & 0 & \dots & 0
\end{pmatrix}
\begin{pmatrix} a_1 \\ \vdots \\ a_n \end{pmatrix}
= \begin{pmatrix} x_1 \\ \vdots \\ x_n \end{pmatrix} \in U,
$$
即 $U = V$，故 $V$ 是单 $M_n(D)$-模.\par
\noindent\fbox{正则模同构于 $nV$} 令
$$
I_i = \left\{ \begin{pmatrix}
0 & \dots & 0 & x_1 & 0 & \dots & 0 \\
\vdots & & \vdots & \vdots & \vdots & & \vdots \\
0 & \dots & 0 & x_n & 0 & \dots & 0
\end{pmatrix}
\in M_n(D) \ \middle| \ x_1, \dots, x_n \in D \right\}.
$$
其中第$i$列非零.则 $I_i$ 是 $M_n(D)$ 的左理想，且有左 $M_n(D)$-模同构 $I_i \cong V, i = 1, \dots, n$.于是有 $M_n(D)$-模同构 $M_n(D) = I_1 \oplus \dots \oplus I_n \cong nV$.\par
\noindent\fbox{$\dim_F V= n[D:F]$} $M_n(D) \cong nV$ 作为 $A$-模，自然也是 $F$-向量空间同构。两边取 $F$-维数：
\[
\dim_F M_n(D) = \dim_F(nV) = n \cdot \dim_F V.
\]
而 $\dim_F M_n(D) = n^2 [D:F]$，故 $\dim_F V = n[D:F]$。\par
\noindent\fbox{$V$唯一} 设 $S$ 是任一单 $A$-模.则
$$
0 \neq S \cong \mathrm{Hom}_A(A, S) \cong \mathrm{Hom}_A(nV, S) \cong n\mathrm{Hom}_A(V, S);
$$
从而 $\mathrm{Hom}_A(V, S) \neq 0$，由 Schur 引理即知 $S \cong V$.这表明 $V$ 是 $A$ 的唯一单模.\par
\noindent\fbox{$\mathrm{End}_A(V) \cong D^{\mathrm{op}}$} 设 $f \in \mathrm{End}_A(V)$.注意到若取$\alpha=\begin{pmatrix}1& 0& \cdots & 0\end{pmatrix}^T$，则$V = A\alpha$.\par
令 $f(\alpha) = \begin{pmatrix} a_1 & \cdots & a_n \end{pmatrix}^T$，$e_{ij}$ 是 $(i, j)$ 处为 $1$、其余全为 $0$ 的 $n$ 阶矩阵.则由
$$
\begin{pmatrix} a_1 \\ \vdots \\ a_n \end{pmatrix}
= f(\alpha) = f(e_{11}\alpha) = e_{11} \begin{pmatrix} a_1 \\ \vdots \\ a_n \end{pmatrix}
= \begin{pmatrix} a_1 \\ 0 \\ \vdots \\ 0 \end{pmatrix},
$$
推出 $a_2 = \dots = a_n = 0$.从而对于任一 $x = \begin{pmatrix} x_1 & \cdots & x_n \end{pmatrix}^T \in V$，有
$$
\begin{aligned}
f(x) &= f((e_{11}x_1 + \dots + e_{n1}x_n)\alpha) = (e_{11}x_1 + \dots + e_{n1}x_n) \begin{pmatrix} a_1 \\ 0 \\ \vdots \\ 0 \end{pmatrix} = \begin{pmatrix} x_1 a_1 \\ \vdots \\ x_n a_1 \end{pmatrix}
= \begin{pmatrix} x_1 \\ \vdots \\ x_n \end{pmatrix} \cdot a_1.
\end{aligned}
$$
由此可见，$f \mapsto a_1$ 给出了 $F$-代数同构：$\mathrm{End}_A(V) \cong D^{\mathrm{op}}$。


\noindent \textbf{(iii)} 将上段证明中的单 $A$-模 $V$ 视为 $A$ 的不可约 $F$-表示，于是得到代数映射 $\rho : A \to \mathrm{End}_F(V)$.\par
\noindent\fbox{$\rho(A) \cong M_n(D)$} 由模 ${}_A V$ 的构造知(因为没有任何非零向量可以使所有列向量变成零)
$$
\mathrm{Ker}\rho = \{ a \in A \mid aV = 0 \} = 0.
$$
因此
$$
\rho(A) \cong A/\mathrm{Ker}\rho = A = M_n(D).
$$
\noindent\fbox{当 $\mathrm{End}_A(V) \cong F$ 时} 由 (ii) 知 $\mathrm{End}_A(V) \cong D^{\mathrm{op}}$，故 $D^{\mathrm{op}} \cong F$，从而 $D = F$，$A = M_n(F)$。此时 $V = F^n$，$\dim_F V = n$。代入上述已证的 $\rho(A) \cong M_n(D)$ 即得 $\rho(A) \cong M_n(F)$——$\rho$ 是 $M_n(F)$ 的忠实表示。

取 $V$ 的 $n$ 个标准单位列向量 $\{e_1,\dots,e_n\}$ 作为 $F$-基。对任意 $a = (a_{ij}) \in M_n(F)$，
\[
\rho(a)(e_j) = a \cdot e_j = \begin{pmatrix} a_{1j} & \cdots & a_{nj} \end{pmatrix}^T,
\]
这正是 $a$ 的第 $j$ 列。因此 $\rho(a)$ 在基 $\{e_1,\dots,e_n\}$ 下的矩阵就是 $a$ 自身。
\end{proof}

\noindent 定理 4.7 的核心信息之一：每个单因子 $M_{n_i}(D_i)$ 上恰有一个单模.既然单模归属于唯一的单因子，那么 $A$ 的互不同构单模总数就等于单因子的个数.由此立即推出——

\begin{cor}
设 $A \cong A_1 \times \dots \times A_n$ 是半单 $F$-代数，其中每个 $A_i$ 均为单代数.则恰有 $n$ 个互不同构的单 $A$-模.

进一步地，若 $\mathrm{Hom}_A(V, V) \cong F$，其中 $V$ 是任一单 $A$-模，则单 $A$-模的个数为 $\dim_F Z(A)$.
\end{cor}

\begin{proof}
\noindent\fbox{单模个数为 $n$} $A \cong A_1 \times \dots \times A_n$，各 $A_i$ 单代数。由推论~4.6，每个 $A_i \cong M_{n_i}(D_i)$，其中 $D_i$ 是可除 $F$-代数。由定理~4.7(ii)，$M_{n_i}(D_i)$ 有唯一的单模——记为 $V_i$。

按定理~4.7(i) 的方式将 $V_i$ 提升为 $A$-模，得到 $n$ 个 $A$-模。这些模互不同构（因 $j \neq i$ 时 $A_j V_i = 0$ 而 $A_j V_j = V_j$，零化子不同）。由定理~4.7(i) 知 $A$ 的任一单模必等同于某个 $V_i$，故恰有 $n$ 个互不同构的单 $A$-模。


\noindent\fbox{分裂条件下单模个数 $= \dim_F Z(A)$} 假设对 $A$ 的每个单模 $V$ 有 $\operatorname{End}_A(V) \cong F$。特别地对上述 $V_i$，$\operatorname{End}_A(V_i) \cong F$。由定理~4.7(ii)，$\operatorname{End}_A(V_i) = \operatorname{End}_{A_i}(V_i) \cong D_i^{\mathrm{op}}$，故 $D_i^{\mathrm{op}} \cong F$，从而 $D_i \cong F$，$A_i \cong M_{n_i}(F)$。

对可除代数的矩阵代数，其中心恰为标量矩阵(3.1习题3)：$Z(M_{n_i}(F)) \cong F$，维数为 $1$。又 $Z(A) \cong Z(A_1) \times \dots \times Z(A_n)$（因为直积中乘法逐分量进行：$(z_1,\dots,z_n)\cdot(a_1,\dots,a_n)=(z_1a_1,\dots,z_na_n)$，故 $z$ 与所有 $a$ 可交换当且仅当每个 $z_i\in Z(A_i)$）。故
\[
\dim_F Z(A) = \sum_{i=1}^n \dim_F Z(A_i) = \sum_{i=1}^n \dim_F Z(M_{n_i}(F)) = \sum_{i=1}^n 1 = n. \qedhere
\]

\end{proof}


\noindent Wedderburn-Artin 定理将半单代数 $A$ 归结为可除代数 $D_1, \dots, D_t$ 和正整数 $n_1, \dots, n_t$，即 $A \cong M_{n_1}(D_1) \times \dots \times M_{n_t}(D_t)$.自然的问题是 $D_i$ 与 $n_i$ 是否由 $A$ 唯一确定？下述定理给出了肯定的回答，从而将半单 $F$-代数的分类归结为有限维可除 $F$-代数的分类.

\begin{thm}
若 $A \cong M_{n_1}(D_1) \times \dots \times M_{n_t}(D_t) \cong M_{m_1}(K_1) \times \dots \times M_{m_s}(K_s)$，其中 $D_1, \dots, D_t$, $K_1, \dots, K_s$ 均是可除 $F$-代数，则 $t = s$，且存在 $\{1, \dots, s\}$ 的置换 $\pi$ 使得 $n_i = m_{\pi(i)}$, $D_i \cong K_{\pi(i)}$.
\end{thm}

\begin{proof}
设 $e_i$ 是 $M_{n_i}(D_i)$ 的单位元，$i = 1, \dots, t$；$f_j$ 是 $M_{m_j}(K_j)$ 的单位元，$j = 1, \dots, s$.则
$$
1 = e_1 + \dots + e_t = f_1 + \dots + f_s.
$$
因为 $e_i, f_j$ 均属于 $A$ 的中心，故 $e_i f_j$ 与 $e_i - e_i f_j$ 均是 $M_{n_i}(D_i)$ 的中心幂等元，且 $e_i f_j (e_i - e_i f_j) = (e_i - e_i f_j) e_i f_j = 0$.于是得到 $M_{n_i}(D_i)$ 的单位元 $e_i$ 的中心幂等元分解 $e_i = e_i f_j + (e_i - e_i f_j)$.\par
从而由命题 1.4 知 $M_{n_i}(D_i)$ 有直积分解，但 $M_{n_i}(D_i)$ 是单代数，故 $e_i f_j = 0$ 或 $e_i f_j = e_i$.类似地，$e_i f_j = 0$ 或 $e_i f_j = f_j$.

因为 $e_i = e_i(f_1 + \dots + f_s)$，故存在 $f_j$ 使得 $e_i f_j \neq 0$.从而 $e_i = e_i f_j = f_j$；若 $l \neq j$，则 $e_i f_l = 0$，这是因为 $e_i f_l = f_j f_l = 0$.这表明对于任一 $i$，存在唯一的 $j$，使得 $e_i = f_j$.于是 $t \leq s$.

同样的讨论可知 $s \leq t$，从而 $t = s$。于是存在 $\{1, \dots, s\}$ 的置换 $\pi$ 使得 $e_i = f_{\pi(i)}$。因此
\[
M_{n_i}(D_i) \cong A e_i = A f_{\pi(i)} \cong M_{m_{\pi(i)}}(K_{\pi(i)}).
\]

至此已得到每对因子的同构。现在只需证：若 $A \cong M_n(D) \cong M_m(K)$，$D$ 与 $K$ 均为 $F$-可除代数，则 $n = m$ 且 $D \cong K$.

根据定理4.7(ii) 知 $A \cong nV \cong mV$，其中 $V$ 是 $A$ 唯一的单模，从而 $n = m$，而且 $D \cong (\mathrm{End}_A(V))^{\mathrm{op}} \cong K$.
\end{proof}

% ========================================== %
% 六、回到群代数
% ========================================== %


\noindent 到这里，半单代数的完整图景已经清晰：结构上就是可除代数上矩阵代数的直积（4.5），单模按单因子分类（4.7），分解还是唯一的（4.10）.现在回到出发点——本节开头提了两个问题：为什么不可约表示数等于共轭类数？为什么维数平方和等于群阶？答案就藏在这个理论中.

最后，将本节的结果应用于有限群的群代数，我们重新得到群表示论中的经典结论.

\begin{thm}
设 $\mathrm{char}F \nmid |G|$.则有限群代数 $FG$ 的互不同构的单模个数 $\leq s$，其中 $s$ 是 $G$ 的共轭类的个数；且若 $F$ 是 $G$ 的分裂域，则 $FG$ 恰有 $s$ 个互不同构的单模.
\end{thm}

\begin{proof}
由条件 $\mathrm{char}F \nmid |G|$ 及 Maschke 定理，群代数 $FG$ 是半单代数。由 Wedderburn-Artin 定理，$FG$ 同构于单代数的直积。设单因子个数为 $n$，则 $FG$ 的互不同构单模个数恰为 $n$（推论~4.8 前半部分）。

推论~4.8 后半部分给出：$n \leq \dim_F Z(FG)$，且当 $F$ 是 $G$ 的分裂域（即 $\operatorname{End}_{FG}(V) \cong F$ 对所有单模 $V$ 成立）时取等号。

\noindent\fbox{$\dim_F Z(FG)=s$} 我们之前证明过一个结论，$Z(FG)$的基为$\{e_i\triangleq \sum_{g\in C_i}g\}$，其中$C_i$是$G$的共轭类，因此$\dim_F Z(FG)=s$.\par
    \begin{proof}
        设$x = \sum_{g \in G} a_g g \in Z(\mathbb{F}G)$，则$\forall\,h \in G$，有$hx = xh$，即$hxh^{-1} = x$.\par
    左式展开得$hxh^{-1} = \sum_{g \in G} a_g (hgh^{-1})$，令$g' = hgh^{-1}$，则$g = h^{-1}g'h$，于是$hxh^{-1} = \sum_{g' \in G} a_{h^{-1}g'h} g'$.\par
    将其与$x = \sum_{g' \in G} a_{g'} g'$比较系数可知，对任意$h, g' \in G$，有$a_{h^{-1}g'h} = a_{g'}$.\par
    这说明如果两个元素相互共轭，它们在中心元素$x$中的系数必定相同.\par
    因此，$x$可以提取出公共系数写为共轭类内部元素和的线性组合，即$x = \sum_i \lambda_i e_i$.\par
    另一方面，易验证对任意共轭类和$e_i$，都有$he_i h^{-1} = e_i$，即$e_i \in Z(\mathbb{F}G)$.\par
    又因为不同的共轭类对应集合互不相交，所以这些类和$\{e_i\}$线性无关，从而构成$Z(\mathbb{F}G)$的基.
    \end{proof}
综上所述$n \leq s, \qquad \text{且 } F \text{ 为分裂域时 } n = s$.
\end{proof}


\noindent 设 $\mathrm{char}F \nmid |G|$ 且 $F$ 是 $G$ 的分裂域.则由 Maschke 定理、Wedderburn-Artin 定理和单模唯一定理知
$$
FG = M_{n_1}(F) \times \dots \times M_{n_s}(F),
$$
其中 $s$ 恰为 $G$ 的共轭类的个数.


\noindent \textbf{收束——回到引子的两个问题：}
\begin{itemize}
    \item \textbf{不可约表示个数}：单因子数 $s$ = 不可约表示数 = $\dim_F Z(FG)$ = 共轭类数.问题 1 得证.
    \item \textbf{维数平方和}：$\dim_F FG = |G| = \sum \dim_F M_{n_i}(F) = \sum n_i^2$.问题 2 得证.
\end{itemize}


令 $e_i$ 是 $M_{n_i}(F)$ 的单位元.则得到 $FG$ 的中心 $Z(FG)$ 的两组基 $\{c_1, \dots, c_s\}$ 和 $\{e_1, \dots, e_s\}$，其中 $c_i$ 是 $G$ 的共轭类 $C_i$ 的类和.这两组基各有优点：$c_1, \dots, c_s$ 有明显的群论意义，而 $e_1, \dots, e_s$ 是正交基.

令
$$
\begin{pmatrix} c_1 \\ \vdots \\ c_s \end{pmatrix}
= \begin{pmatrix}
a_{11} & \dots & a_{1s} \\
\dots & \dots & \dots \\
a_{s1} & \dots & a_{ss}
\end{pmatrix}
\begin{pmatrix} e_1 \\ \vdots \\ e_s \end{pmatrix}, \quad
\begin{pmatrix} e_1 \\ \vdots \\ e_s \end{pmatrix}
= \begin{pmatrix}
b_{11} & \dots & b_{1s} \\
\dots & \dots & \dots \\
b_{s1} & \dots & b_{ss}
\end{pmatrix}
\begin{pmatrix} c_1 \\ \vdots \\ c_s \end{pmatrix},
$$
则 $A = (a_{ij})_{s \times s}$ 和 $B = (b_{ij})_{s \times s}$ 是 $F$ 上互逆的 $s$ 阶矩阵.下面的命题给出了 $A, B$ 与 $G$ 的特征标表的关系.

\begin{prop}
设 $h_i = |C_i|$.则
$$
a_{ij} = \frac{h_i}{n_j} \chi_j(g_i), \quad b_{ij} = \frac{n_i}{|G|} \chi_i(g_j^{-1}),
$$
其中 $\chi_i$ 是单因子 $M_{n_i}(F)$ 上的单模 $V_i$（自然地视为 $FG$-模）的特征标，$g_i$ 是 $C_i$ 中任一元.
\end{prop}

\begin{proof}
若 $j \neq k$，由定理4.7(i)$M_{n_j}(F)$ 在 $V_k$ 上的作用为零，而 $e_j$ 在 $V_j$ 上的作用是恒等作用.因此
$$
\chi_j(e_k) := \mathrm{tr}(e_k) = \begin{cases}
n_j, & k = j, \\
0, & k \neq j.
\end{cases}
$$
从而
\[
\chi_j(c_i) = \chi_j \Bigl( \sum_{k=1}^s a_{ik}e_k \Bigr) = \sum_{k=1}^s a_{ik}\, \chi_j(e_k) = a_{ij}n_j. \tag{1}
\]

另一方面，$c_i = \sum_{g \in C_i} g$ 是共轭类 $C_i$ 中 $h_i = |C_i|$ 个元素的和，对有限维表示的特征标有可加性：
\[
\chi_j(c_i) = \chi_j\Bigl(\sum_{g \in C_i} g\Bigr) = \sum_{g \in C_i} \chi_j(g) = h_i\, \chi_j(g_i), \tag{2}
\]
其中 $g_i$ 为 $C_i$ 中任一元（共轭元具有相同的特征标）。

比较 (1) 和 (2) 得 $a_{ij} n_j = h_i \chi_j(g_i)$。由引理~II.3.3.A，$\mathrm{char}F \nmid n_j$，故 $n_j$ 在 $F$ 中可逆，于是
\[
a_{ij} = \frac{h_i}{n_j}\,\chi_j(g_i).
\]

\noindent 下面求 $b_{ij}$。两路计算 $\chi_{\mathrm{reg}}(e_i g_j^{-1})$。

\smallskip
\noindent\fbox{用 $c_k$ 基展开 $e_i$} 由 $e_i = \sum_{k=1}^s b_{ik} c_k$，得
\[
e_i g_j^{-1} = \sum_{k=1}^s b_{ik}\, c_k g_j^{-1}.
\]
$c_k = \sum_{g \in C_k} g$，故 $c_k g_j^{-1} = \sum_{g \in C_k} g g_j^{-1}$。正则特征标满足
\[
\chi_{\mathrm{reg}}(g) = \begin{cases} |G|, & g = 1, \\ 0, & g \neq 1. \end{cases}
\]
因此 $\chi_{\mathrm{reg}}(g g_j^{-1}) \neq 0$ 当且仅当 $g g_j^{-1} = 1$，即 $g = g_j$。这当且仅当 $g_j \in C_k$ 时出现，且此时该项贡献 $\chi_{\mathrm{reg}}(1) = |G|$，其余 $g \neq g_j$ 的项贡献 $0$，故 $\chi_{\mathrm{reg}}(c_k g_j^{-1}) = |G|$。由于 $g_j$ 只属于一个共轭类，满足 $g_j \in C_k$ 的 $k$ 是唯一的——记这个 $k$ 恰为 $j$。于是
\[
\chi_{\mathrm{reg}}(e_i g_j^{-1})
= \sum_{k=1}^s b_{ik}\, \chi_{\mathrm{reg}}(c_k g_j^{-1})
= b_{ij} \cdot |G| + \sum_{k \neq j} b_{ik} \cdot 0
= |G|\, b_{ij}. \tag{3}
\]

\smallskip
\noindent\fbox{用 $\chi_{\mathrm{reg}} = \sum_{k=1}^s n_k \chi_k$ 分解} 因 $e_i g_j^{-1} \in M_{n_i}(F)$，而$M_{n_i}(F)$ 只作用在 $V_i$ 上非零，故对 $k \neq i$ 有 $\chi_k(e_i g_j^{-1}) = 0$。对 $k = i$，$e_i$ 在 $V_i$ 上的作用为恒等，故
\[
(e_i g_j^{-1}) \cdot V_i = g_j^{-1} \cdot (e_i \cdot V_i) = g_j^{-1} \cdot V_i,
\]
即 $\chi_i(e_i g_j^{-1}) = \chi_i(g_j^{-1})$。于是
\[
\chi_{\mathrm{reg}}(e_i g_j^{-1})
= \sum_{k=1}^s n_k \chi_k(e_i g_j^{-1})
= n_i \chi_i(e_i g_j^{-1})
= n_i \chi_i(g_j^{-1}). \tag{4}
\]

\smallskip
\noindent 比较 (3) 和 (4) 得 $|G|\, b_{ij} = n_i \chi_i(g_j^{-1})$，故
\begin{equation*}
b_{ij} = \frac{n_i}{|G|}\,\chi_i(g_j^{-1}). \qedhere
\end{equation*}
\end{proof}

引理II.3.3.A的证明
\begin{proof}
    由于$\rho_f$是$V\to V$的G-模同态 且$V$不可约，由Schur引理，$\rho_f = \lambda_f\cdot 1_V$.\par
    则$\begin{aligned}[t]
         n\cdot \lambda_f &= \mathrm{tr}\rho_f = \mathrm{tr}\left(\frac{1}{|G|}\sum_{g\in G}f(g)\rho(g)\right) = \frac{1}{|G|}\sum_{g\in G}f(g)\mathrm{tr}\rho(g) \\
         &= \frac{1}{|G|}\sum_{g\in G}f(g)\chi(g) = \frac{1}{|G|}\sum_{g\in G}f(g)\chi^*(g^{-1})\\ 
         &= (f,\chi^*)
    \end{aligned}$.\par
    $\Rightarrow \lambda_f = \frac{(f,\chi^*)}{n}$.\par
    取$f=\chi^*$，则$n\lambda_f = (\chi^*,\chi^*) = 1$，从而$\mathrm{char}\,\mathbb{F}\nmid n$.
\end{proof}

% ========================================== %
% 七、推广与展望
% ========================================== %

\bigskip
\noindent \textbf{审视证明中真正用到的条件.} 回顾 (i) $\Rightarrow$ (ii) $\Rightarrow$ (iii) 的证明，实际上只用到了：

\begin{enumerate}
    \item \textbf{降链条件 (DCC)} —— 即 $A$ 有 Artin 性（保证极小左理想存在、合成列存在）；
    \item \textbf{左正则模是半单模} —— 即 $A$ 没有非平凡的 Jacobson 根（$\mathrm{rad}(A) = 0$）.
\end{enumerate}

有限维性只是 Artin 性 + Noether 性的充分条件，而非必要条件.这暗示定理有更广的适用范围.

\medskip
\noindent \textbf{Wedderburn-Artin 定理的完全形态}（Artin, 1927）：

\begin{thm*}
设 $R$ 是环.则下述等价：
\begin{enumerate}
    \item $R$ 是半单环（左 Artin 环且任意左 $R$-模是半单模）；
    \item $R$ 是左 Artin 环且 Jacobson 根 $\mathrm{rad}(R) = 0$；
    \item $R \cong M_{n_1}(D_1) \times \dots \times M_{n_t}(D_t)$，其中 $D_i$ 为可除环.
\end{enumerate}
\end{thm*}

特别地，对于任意左 Artin 环 $R$，$R/J(R)$ 总是半单环，因此同构于可除环上矩阵环的直积.



\end{document}
